\section{Auslegung der Steuerung und Überwachung eines Hochleistungs-Gleichrichters}

\subsection{Aufgabenstellung und Ausgangslage}
Zum Einstieg in das CAE-System \textit{EPLAN Electric P8 (Version 2025)} und die betrieblichen Standards wurde der Schaltplan für die Steuerung und Überwachung eines Hochleistungs-Gleichrichters erstellt.
Als Grundlage diente ein Schaltungsentwurf des Auftraggebers.

Die Aufgabe bestand darin, diesen Entwurf auf fertigungstechnische Machbarkeit zu prüfen, Korrekturen vorzunehmen, einen vollständigen Schaltplan zu erstellen und die passenden Betriebsmittel auszuwählen.

\subsection{Normgerechte Strukturierung nach DIN EN 81346--2}
Der Schaltplan wurde nach der Norm \textbf{DIN EN 81346--2} aufgebaut.
Dadurch wurden die Betriebsmittel anhand ihrer Funktion und Einbauortes eindeutig gekennzeichnet.
Dies wurde so gemacht, da dadurch die Fertigung und Inbetriebnahme vereinfacht wird und die Dokumentation für den späteren Betrieb und die Wartung des Gleichrichters normgerecht ist.

\subsection{Schaltungsprüfung und konstruktive Anpassungen}
Der ursprüngliche Entwurf enthielt einige Details, die in der Fertigung zu Problemen geführt hätten. Folgende Änderungen wurden umgesetzt:
\begin{itemize}
    \item \textbf{Klemmenanschlüsse bereinigen:} Im Entwurf waren teilweise zwei Adern auf einen Klemmenanschluss gelegt.
          Da Reihenklemmen meist nur für einen Leiter pro Klemmstelle zugelassen sind, wurden zusätzliche Klemmen eingeplant.
          Um Kosten zu sparen, wurde dort, wo es möglich war, auf die Verwendung von Steg- und Aderbrücken gesetzt.
    \item \textbf{Trennung der Stromkreise:} Um Verdrahtungsfehler bei der Montage zu verhindern und Störeinflüsse zu vermeiden, wurden die Klemmenleisten nach Spannung und Funktion getrennt aufgebaut:
          \begin{itemize}
              \item \SI{400}{\volt} AC (Leistungsteil)
              \item \SI{24}{\volt} DC (Steuerspannung)
              \item Mess- und Meldesignale
          \end{itemize}
\end{itemize}

\subsection{Auswahl und Dimensionierung der Komponenten}
Auf Basis der technischen Daten des Gleichrichters und der Steuerungskomponenten wurden die Betriebsmittel ausgewählt und überprüft:
\begin{itemize}
    \item \textbf{Schalt- und Schutzgeräte:} Schütze wurden passend zur geschalteten Last und der jeweiligen Gebrauchskategorie (AC-1 bzw. AC-3) dimensioniert.
          Leitungsschutzschalter und Sicherungen wurden bezüglich Nennstrom und Auslösecharakteristik gewählt.
    \item \textbf{Überspannungsschutz:} Für den weiteren Schutz wurden Überspannungsableiter eingeplant.
    \item \textbf{Leitungen und Klemmen:} Querschnitte von Leitungen und Klemmen wurden anhand der Strombelastbarkeit und der Verlegeart gewählt.
          Für Signalleitungen wurden geschirmte Typen vorgesehen.
\end{itemize}

\subsection{Artikelverwaltung und Anbindung an FactWork}
Damit der Schaltschrank wirtschaftlich gefertigt werden kann, wurden überwiegend Bauteile ausgewählt, die bei F.EE bereits als Standardartikel hinterlegt sind.

Die Artikeldaten wurden über die Schnittstelle zum internen ERP-System \textit{FactWork} in EPLAN übernommen:
\begin{itemize}
    \item Alle relevanten Daten (Bestellnummern, Hersteller, technische Kennwerte) sind direkt mit den Schaltplan-Symbolen verknüpft.
    \item Nach Fertigstellung des Plans können Stücklisten, Klemmenpläne und Kabelpläne automatisch generiert und direkt an den Einkauf und den Schaltschrankbau übergeben werden.
\end{itemize}